\chapter{Discussion and Conclusions}\label{sec:discussion-and-conclusions}

\paragraph{Scope of the discussion.} The conclusions discussed in this chapter draw on three parts of the thesis. Chapter~\ref{sec:experiments}'s canonical comparison evaluates relay strategies under the canonical memoryless channel model, where transmission is uncoded, the noise is white, and processing is performed on a per-symbol basis. The same chapter's supplementary coded study (Section~\ref{sec:coded-block-df}) departs from that model by introducing a rate-1/2 convolutional code, a block-DF relay, and an adaptive modulation-and-coding scheme. Chapter~\ref{sec:unknown-channels} departs from it separately by introducing impairments such as intersymbol interference (ISI), nonlinear distortion, and unknown phase. Accordingly, statements concerning channel memory, sequence detection, Viterbi MLSE, pilot-based estimation, or blind equalization should be understood as applying only to the extension study of Chapter~\ref{sec:unknown-channels}; statements concerning coding, MCS adaptation, or latency budgets apply only to the coded study; neither applies to the canonical memoryless setup.

The canonical core (Chapter~\ref{sec:experiments}) settles hypotheses H1--H4 and establishes the baseline finding that, on a matched known channel, a minimal learned relay does not improve on classical decode-and-forward. The thesis's principal contribution (Chapter~\ref{sec:unknown-channels}) settles H5 and is where the substantive new claim lies: a systematic delineation of when a learned relay adds value under unknown and mismatched channels. The interpretation below treats these two as the thesis's two principal contributions, with the unknown-channel results carrying the main argument; the supplementary coded study is the smaller, third part, discussed separately below where it bears on the relay-decoding and rate-adaptation questions it was designed to measure.

\section{Interpretation of Results}\label{sec:interpretation-of-results}

The experimental results reveal several consistent patterns across the canonical setup, its supplementary coded-link study, and the unknown-channel contribution. This section interprets these patterns through the lens of the theoretical framework established in Section \ref{sec:theoretical-foundations-channel-models}  and evaluates each research hypothesis.

\subsection{Low-SNR Advantage of Neural Relays over AF (H1: Confirmed)}\label{sec:low-snr-advantage-of-neural-relays-h1-confirmed}

\textbf{Low SNR (0-4 dB): AI advantage over AF.} At low SNR, AI-based relays consistently outperform AF on the canonical channel, confirming H1. The mechanism is the structure of the MMSE-optimal (posterior-mean) relay estimate. The scalar AWGN form below is \emph{qualitative motivation}, not the exact canonical estimator: it holds for a fixed noise variance, whereas the canonical link is Rayleigh-faded, so the exact posterior is conditioned on the retained fading magnitude --- $\mathbb{E}[x \mid \tilde{y}, |h|{=}a] = \tanh(a^{2}\tilde{y}/\sigma^{2})$ under the zero-forcing convention of Section~\ref{sec:system-model-1}, where $\tilde{y} = y/h$ carries conditional noise variance $\sigma^{2}/a^{2}$ rather than $\sigma^{2}$ --- or, where that magnitude is not supplied to the relay explicitly, marginalizes over it to give a fading-mixture posterior rather than a single fixed-variance $\tanh$. What the trained network approximates empirically is that mixture. With that scope stated, the fixed-variance form makes the shape of the mapping plain:

\begin{equation}\protect\phantomsection\label{eq:optimal-denoiser-relay}{f^*(y) = \mathbb{E}[x | y] = \tanh\left(\frac{y}{\sigma^2}\right)
}\end{equation}
\addcontentsline{equ}{listequations}{\protect\numberline{\theequation}\quad optimal-denoiser-relay}

At low SNR (e.g., 0 dB, \(\sigma^2 = 1\)), this is a gentle sigmoid: \(f^*(y) = \tanh(y)\). The fading case rescales its slope per symbol and, unconditioned, softens it into a mixture, but the qualitative point is unchanged: the optimal relay map is a smooth saturating non-linearity rather than a linear gain. The neural network can approximate such a mapping accurately, producing a \textbf{soft estimate} that preserves information about the confidence of the decision. By contrast:

\begin{itemize}
\item
  \textbf{DF} applies the hard-decision function \(f_{\text{DF}}(y) = \text{sign}(y)\), which discards all soft information. At low SNR, many received samples lie near the decision boundary (\(|y| \approx 0\)), and DF assigns them full confidence (\(\pm 1\)) regardless of the actual uncertainty. This information loss leads to excess errors on the second hop.
\item
  \textbf{AF} applies \(f_{\text{AF}}(y) = Gy\), which is linear and preserves the noisy signal structure but amplifies the noise by the same factor as the signal. The effective SNR degradation is given by \(\text{SNR}_{\text{eff}} = \gamma^2 / (2\gamma + 1) \approx \gamma/2\) at high SNR: a 3 dB penalty.
\end{itemize}

The AI relay implements a \textbf{non-linear soft mapping} that is intermediate between these extremes: it suppresses noise (like DF's regeneration) while preserving soft information near the decision boundary (unlike DF's hard decision). This explains why learned non-linear relay processing can outperform AF and, under selected channel conditions, also outperform DF at low SNR.

Among the AI methods, the low-SNR advantage over AF is realized by the \emph{minimal} feedforward relays, not by the larger sequence models. On the canonical Rayleigh channel (Table~\ref{tbl:table2}, 0~dB) the Hybrid, MLP and VAE relays reach $0.3329$, $0.3349$ and $0.3359$ respectively, essentially matching symbol-wise DF's $0.3337$ and improving substantially on AF's $0.4076$, while the higher-capacity Transformer, Mamba-S6 and Mamba-2 (SSD) relays sit at $0.4014$, $0.4071$ and $0.4032$ -- that is, at essentially AF's level, capturing almost none of the available advantage. The advantage over AF is therefore \emph{not} driven by parameter count: the smallest network already captures it, consistent with the complexity ordering discussed under H3 and with the convergence observed in the equal-parameter-budget comparison (H4, Section~\ref{sec:normalized-parameter-comparison}).

\subsection{Classical Dominance at High SNR (H2: Confirmed)}\label{sec:classical-dominance-at-high-snr-h2-confirmed}

\textbf{Medium-to-high SNR (≥6 dB): Classical dominance.} At medium and high SNR, DF matches or exceeds all AI methods with exactly zero trainable parameters and zero training time. The theoretical explanation is straightforward: at high SNR, \(\sigma^2 \to 0\) and the posterior-mean estimate on each real decision axis tends to the corresponding hard slice. The symbol-wise DF relay therefore implements the limiting \emph{hard symbol-regenerating} rule in this regime, while any neural network approximation introduces a small but non-zero reconstruction error due to the finite precision of the learned mapping and the tanh output saturation (which approaches but never reaches \(\pm 1\)).

The effect is visible analytically in the AWGN limit of the model ($|h|\equiv1$), where at 10 dB the two-hop DF BER is \(2Q(\sqrt{10}) \cdot (1 - Q(\sqrt{10})) \approx 0.0016\), and empirically on the canonical Rayleigh channel, where DF attains the lowest or tied-lowest BER at every point from 6 dB upward. The AI models cannot beat DF in this regime for the uncoded matched setting because componentwise hard slicing is the MAP symbol decision, and thus the optimal symbol-regenerating DF rule, after coherent equalization (Remark~\ref{rem:df-terminology}); they can only match it, and the larger sequence models slightly underperform due to the increased variance of their more complex learned mappings.

\subsection{Robustness Across the Fading Ensemble}\label{sec:channel-robustness}

\textbf{Robustness across fades.} The relay ranking on the canonical channel is stable across the full evaluated SNR range, even though each transmitted symbol experiences an independent fading realization. The mechanism is that after coherent compensation (\(\hat{x} = y \cdot h^*/|h|\)), the relay processes a signal with AWGN-like noise at a \emph{random effective SNR}; the network, trained across multiple SNR values, has learned a denoising function that adapts to different noise levels, so the Rayleigh ensemble (a mixture of effective SNRs weighted by the fading distribution) is handled by virtue of multi-SNR training alone, with no per-realization retraining or explicit CSI input.

\section{Capacity and Overfitting (H3: Not Supported)}\label{sec:capacity-and-overfitting}

One of the most practically useful findings is that relay BER does not improve with model size over the range evaluated: the Minimal MLP architecture (169 parameters) matches the performance of models with 100--140$\times$ more parameters (3K--24K), and no larger evaluated model improved on it under the training protocol used. The stronger reading of H3 --- that the curve turns upward because excess capacity causes \emph{overfitting} --- is not established by the evidence collected here. Each model was trained once, with no seed replication, no train/test-gap measurement and no held-out validation curve, so a size effect cannot be separated from run-to-run variance or from an optimization difficulty that more capacity happens to introduce. The claim carried forward is therefore the weaker, directly measured one: larger evaluated models did not improve BER. What follows in this section is offered as a candidate explanation for that observation, not as a demonstrated mechanism.

This result has important theoretical and practical implications:

\subsection{Information-Theoretic Perspective}\label{sec:information-theoretic-perspective}

The canonical memoryless relay-denoising task maps a window of \(2w+1\) real-valued noisy observations to a single real-valued clean estimate, applied per axis to the canonical QPSK symbol (equivalently, the real BPSK-over-AWGN form of the same problem). Each real axis carries exactly 1 bit of information, and under the canonical i.i.d.\ model the sufficient statistic for the current bit is the current observation alone, with Bayes-optimal posterior mean \(f^*(y_i) = \tanh(y_i / \sigma^2)\). The neighboring window samples used by the learned relays are therefore redundant architectural inputs in this setting rather than part of the sufficient statistic, so the \textbf{effective dimensionality} of the target mapping is still 1: far below the \(2w+1 = 5\) or \(11\) input dimensions.

As a heuristic, the minimum description length (MDL) principle suggests matching model complexity to the effective dimensionality of the target mapping. Since the target here is a one-dimensional monotone function of the current sufficient statistic, even the 169-parameter network has capacity far in excess of what the mapping requires; this is offered as an intuition for why the larger models did not improve on the smallest one, not as a formal MDL bound.

\subsection{Bias-Variance Analysis}\label{sec:bias-variance-analysis}

The bias-variance decomposition \cite{GoodfellowBengioCourville2016DeepLearning} provides a quantitative framework:

\begin{equation}\protect\phantomsection\label{eq:bias-variance-relay}{\text{MSE}(\hat{x}) = \underbrace{(\mathbb{E}[\hat{x}] - f^*(y))^2}_{\text{Bias}^2} + \underbrace{\mathbb{E}[(\hat{x} - \mathbb{E}[\hat{x}])^2]}_{\text{Variance}} + \underbrace{\sigma^2_{\text{irred}}}_{\text{Noise floor}}
}\end{equation}
\addcontentsline{equ}{listequations}{\protect\numberline{\theequation}\quad bias-variance-relay}

\begin{itemize}
\tightlist
\item
  \textbf{169 parameters:} Low bias (sufficient capacity for the simple \(f^*\)) and low variance (limited capacity prevents memorization). The model operates at the optimal bias-variance trade-off.
\item
  \textbf{3,000 parameters:} Similar bias (the target function hasn't changed) and slightly higher variance. Performance is nearly identical to 169 params.
\item
  \textbf{24,001 parameters (Mamba-S6):} Despite having 140$\times$ more parameters than the MLP, it does not improve on it; on the canonical channel it is worse at low SNR (Table~\ref{tbl:table2}). The marginal utility of the additional parameters is at best vanishingly small.
\end{itemize}

This decomposition is a framework for interpreting the observation, not a measurement of it: bias and variance were not estimated separately, which would require training each configuration repeatedly across seeds and comparing training against held-out error. An earlier version of this section additionally attributed ``clear overfitting'' to an 11{,}201-parameter Maximum MLP; that configuration is not part of the reported experiments and no such measurement exists, so the claim has been withdrawn rather than restated.

\subsection{Does the Complexity Curve Turn Upward?}\label{sec:complexity-curve-turn}

The preceding subsections offer bias-variance and MDL readings of H3 but note that the stronger, overfitting reading --- BER rising again once capacity exceeds what the mapping requires --- was not established, because each model had been trained once and a size effect could not be separated from run-to-run variance. A dedicated size sweep was therefore run to supply the missing evidence. It covers the nine channel families used across Chapters~\ref{sec:experiments} and~\ref{sec:unknown-channels} plus the coded scenario, varying relay window, depth and width independently, with three independent initializations at every configuration and each configuration held to its \emph{worst} initialization, so that a size counts only if it works whichever way it initializes --- a necessary precaution given the spread reported below. Degradation is measured as the additional SNR required to reach a given BER, which unlike a relative BER penalty does not divide by a vanishing denominator at high SNR.

Because seed replication is the specific evidence this section previously lacked, the across-initialization spread is treated as a measurement in its own right rather than as a caveat. Table~\ref{tbl:seed-spread}, in Appendix~\ref{sec:appendix-f-minimum-size}, gives per channel, the spread in SNR penalty between the best and worst of the three initializations at the same configuration, summarized over all configurations swept.

Two things follow. First, initialization variance is not a small correction: across all 152 configurations the median spread is \(0.180\)~dB and \(57\%\) exceed \(0.15\)~dB, with a worst case of \(3.18\)~dB on the complex ISI channel at eleven parameters. On the nonlinear-bias channel the spread is \(1.23\)~dB at essentially every configuration, which is the two-attractor behaviour noted below rather than a gradual variance. Second, the variance is strongly channel-dependent: the canonical Rayleigh channel is highly reproducible at \(0.008\)~dB median, and the composite and AWGN channels nearly so, while four channels sit above \(0.3\)~dB. Any single-run size comparison on the latter group measures which initialization was drawn, not which architecture is better --- which is precisely the confound this section identified in the original H3 evidence, now quantified.

Against that yardstick, \textbf{the upward arm is not observed}, and the sweep had the resolution to see it had it been there. Holding the window fixed so that capacity is the only variable, the worst-case penalty against the classical baseline is flat on the memoryless channels and monotonically decreasing on the channels with memory. On the canonical Rayleigh channel at window~5 the penalty moves from \(+0.03\)~dB at one hidden unit to \(+0.06\)~dB at twenty-four; on the composite channel at window~7 it improves monotonically from \(-1.97\)~dB to \(-2.23\)~dB as width grows from four to forty-eight. Where a rising tail appears at all it is at most \(0.15\)~dB, which is below the \(0.180\)~dB median across-initialization spread of Table~\ref{tbl:seed-spread} and therefore not separable from the draw. The one clearly non-monotone case, the nonlinear-bias channel, oscillates between two distinct outcomes (\(-1.63\) and \(-2.87\)~dB) rather than tracing a curve, which is the signature of training landing in one of two attractors rather than of capacity being harmful.

Two corrections to the reading of H3 follow. First, the weaker claim is now confirmed under seed replication rather than merely unrefuted: additional parameters neither help nor hurt on a memoryless channel over roughly two orders of magnitude of model size. Second, an effect that appears substantial under a relative BER criterion largely dissolves under the SNR-penalty criterion. On AWGN the penalty grows from \(+11\%\) to \(+28\%\) in relative BER as the window widens, which invites the conclusion that a window is actively harmful; the same measurements correspond to \(+0.06\)~dB and \(+0.11\)~dB. The effect is real in sign but negligible in magnitude, and the relative criterion overstates it because the absolute BER gap it divides by falls to \(4\times10^{-4}\).

What the sweep does establish is that the binding resource is not parameter count but \textbf{window matched to channel memory}. On the memoryless channels the smallest configuration examined --- four parameters, a single hidden unit with no window --- is within \(0.15\)~dB of the classical baseline on three of five channels and within \(0.03\)~dB on the canonical channel, some \(42\times\) smaller than the 169-parameter relay of Table~\ref{tbl:table2}. On the three-tap channels the same four-parameter configuration is between \(11\) and \(13\)~dB adrift, and closing that gap requires widening the window rather than the layer: on the real ISI channel the penalty falls from \(+12.95\)~dB at window~1 to \(+1.67\)~dB at window~7. Depth was varied from one to three hidden layers and did not improve any configuration; where a two-layer network matched, a single layer matched at fewer parameters.

The coded scenario behaves differently again and is the one case where size is demonstrably not the operative variable. Against block DF, every configuration from 18 to 1{,}076 parameters lies between \(+1.86\) and \(+1.92\)~dB, so a sixty-fold range in model size buys \(0.06\)~dB. The relay is not capacity-limited there but interface-limited: it commits to a hard constellation decision per symbol, at \(90\)--\(94\%\) symbol accuracy, before the outer code is decoded, so the destination decoder receives no reliability information that block DF would have exploited. This identifies a design change rather than a scaling law, and it is taken up in Section~\ref{sec:future-work}.

Taken together, these measurements support the weaker claim and leave the overfitting hypothesis unsupported --- now on the basis of a sweep designed to detect it, rather than for want of evidence.

The framing that replaces it is not that smaller models are better, but that \textbf{architectural complexity is only exercised by channel memory}. Capacity is neither good nor bad in itself; whether it can be used at all depends on the channel. On a memoryless channel the sufficient statistic for the current symbol is the current observation, so there is nothing for additional window taps to represent and nothing for additional units to fit: every configuration from four parameters to several hundred lands within a fifth of a decibel of the same value, and the extra capacity is not harmful so much as idle. Introduce three taps of intersymbol interference and the same capacity becomes load-bearing --- the relay must now represent a function of several consecutive samples, a window-1 network cannot express it at any width, and widening the window is worth between $8$ and $11$~dB. The size axis looks flat in the first case and decisive in the second because the channel, not the architecture, determines whether the parameters have work to do.

This reading accounts for both halves of the sweep, which ``less is more'' does not: a principle that smaller is better would predict the memoryless behaviour and mispredict the ISI channels, where the minimum rises from four parameters to between $73$ and $145$ and the wider network is unambiguously the better one. It also explains the shape of the failure at the two extremes. Where the window is too narrow for the channel's memory the relay fails structurally rather than gradually, as on the unknown-phase channel of Section~\ref{sec:flat-unknown-channels}, whose information symbol is a product of two consecutive samples that a one-tap network cannot represent at any width. Where the window is far wider than the memory, as at window 31 on the coded channel, the additional inputs carry no information about the current symbol and the optimization becomes harder rather than the model more expressive. The design rule is therefore to match the window to the channel's impulse response first and to treat width as a secondary knob, which is a statement about the channel as much as about the network.

\subsection{Practical Implications}\label{sec:practical-implications}

\begin{enumerate}
\def\labelenumi{\arabic{enumi}.}
\item
  \textbf{Occam's Razor for relay AI.} The canonical relay denoising task has inherently low complexity: the per-axis mapping from noisy to clean symbols is fundamentally simple, and the neural network needs only enough capacity to learn a non-linear threshold function with some contextual smoothing. Adding parameters beyond this minimal requirement bought no measurable accuracy anywhere in this thesis.
\item
  \textbf{Deployment feasibility.} A 169-parameter model requires approximately 0.7 KB of memory (169 float32 values) and trains in seconds on a CPU. This makes AI-based relay processing viable even on severely resource-constrained embedded relay nodes, such as IoT devices or sensor network relays. The model can be stored in on-chip SRAM and executed without any external memory access.
\item
  \textbf{Generalization.} The 169-parameter MLP handles the full Rayleigh fading ensemble, every effective SNR the fades produce, with a single set of weights and no per-condition retraining, an instance of generalization that would be harder to achieve with a larger, more specialized model.
\end{enumerate}

\section{Practical Deployment Recommendations}\label{sec:practical-deployment-recommendations}

Based on the comprehensive evaluation, we propose the following deployment strategy:

\begin{longtable}[]{@{}
  >{\raggedright\arraybackslash}p{(\linewidth - 4\tabcolsep) * \real{0.3333}}
  >{\raggedright\arraybackslash}p{(\linewidth - 4\tabcolsep) * \real{0.3333}}
  >{\raggedright\arraybackslash}p{(\linewidth - 4\tabcolsep) * \real{0.3333}}@{}}
\caption[Practical deployment recommendations]{Practical deployment recommendations: recommended relay strategy by operating regime}\label{tbl:table32}\tabularnewline
\toprule\noalign{}
\begin{minipage}[b]{\linewidth}\raggedright
Operating Regime
\end{minipage} & \begin{minipage}[b]{\linewidth}\raggedright
Recommended Relay
\end{minipage} & \begin{minipage}[b]{\linewidth}\raggedright
Rationale
\end{minipage} \\
\midrule\noalign{}
\endfirsthead
\toprule\noalign{}
\begin{minipage}[b]{\linewidth}\raggedright
Operating Regime
\end{minipage} & \begin{minipage}[b]{\linewidth}\raggedright
Recommended Relay
\end{minipage} & \begin{minipage}[b]{\linewidth}\raggedright
Rationale
\end{minipage} \\
\midrule\noalign{}
\endhead
\bottomrule\noalign{}
\endlastfoot
\textbf{Low SNR (0-4 dB)} & MLP Minimal / Hybrid / channel-specific selection & Neural relays often help; exact best choice is channel-dependent, with DF remaining strong on some fading scenarios \\
\textbf{Medium SNR (4-8 dB)} & Hybrid (MLP + DF) & Automatic switching at empirically chosen crossover\\
\textbf{High SNR (\textgreater8 dB)} & DF & Zero parameters; lowest BER among the evaluated per-symbol relays \\
\textbf{Resource-constrained} & MLP Minimal (169 params) & 0.7 KB memory, \textless5 s training \\
\textbf{Best overall (BPSK/QPSK)} & Hybrid & Combines AI advantage with DF reliability \\
\end{longtable}

The Hybrid relay is recommended as the default deployment choice, where the operating SNR can be estimated reliably, because it automatically selects MLP processing at low SNR and DF at high SNR, achieving near-optimal performance across the entire SNR range with minimal complexity.

\subsection{Latency and Adaptivity Considerations}\label{sec:latency-and-adaptivity-considerations}

Two practical questions bear on whether a learned relay can be deployed with minimal impact: the latency it adds relative to AF/DF, and its ability to adapt to varying channel conditions without retraining.

\textbf{Latency.} Two distinct latency components must be separated. The first is \emph{buffering} latency, incurred by the sliding window $y_R[i-w:i+w]$. In the half-duplex protocol studied here this component is absorbed by the protocol itself: the relay buffers the entire listening-phase block before the transmission phase begins (Remark~\ref{rem:window-causality}), so the window's $w$-symbol look-ahead adds no delay beyond what store-and-forward operation already imposes. Only in a streaming or ultra-low-latency protocol, outside the scope of this thesis, would the symmetric window cost $w$ symbols of look-ahead that the memoryless AF/DF baselines ($w=0$) do not pay. The second component is \emph{compute} latency: inference cost per relayed symbol. Here the minimal-capacity result (H3) is decisive: since the 169-parameter MLP matches the performance of models $100-140\times$ larger, the deployed network reduces to two small matrix--vector products per symbol, a per-symbol cost of a few hundred multiply-accumulates, small in absolute terms even if it cannot literally match the single comparison of DF's $\operatorname{sign}(y)$. The compute-latency argument therefore holds \emph{because of}, not despite, the less-is-more finding; it would not hold for the 24K-parameter sequence models, whose recurrent or attention-based structure imposes materially higher per-symbol cost for no BER benefit at the relay window length.

The coded study of Section~\ref{sec:coded-latency-throughput} adds a third component that the canonical comparison never had to consider: \emph{buffering}. A block relay cannot emit until it has decoded a whole frame ($202$ symbols there), and the soft BCJR variant cannot even in principle, its backward recursion running from the end of the frame. A windowed learned relay needs only its $w$-symbol look-ahead, ten symbols, and that figure is constant in frame length where the block relays' grows with it.

Section~\ref{sec:joint-latency-memory} measures these components together on a channel with memory, and two of the arguments above do not survive it. Treating buffering as the learned relay's advantage presumes that sequence detection is latency-expensive, and it is not: bounded-traceback MLSE reaches full accuracy after three symbols of decision delay, fewer than the window-11 relay's five, so a latency budget tight enough to exclude MLSE excludes the learned relay first. Latency separates the block relays from everything else, and nothing else from anything. The appeal to wall-clock measurement fares no better --- that the measured advantage over Viterbi spans a factor of $50$ to $63$ \emph{depending on the machine} (Table~\ref{tbl:table41}) is itself the objection --- and counting multiply-accumulates instead turns the argument into a threshold at $L^{\star}=3.35$ taps (Equation~\ref{eq:mac-crossover}), below which the learned relay is worse on cost, delay and accuracy at once. The measurement and its two tables are given in Section~\ref{sec:joint-latency-memory} and not repeated here; what matters for deployment is that the learned relay's cost advantage is real but conditional, and the condition is a property of the channel rather than of the processor it runs on.

\textbf{Adaptivity.} A single network, trained across the range of simulated SNRs, handles the entire canonical fading ensemble without retraining or explicit CSI input (Section~\ref{sec:channel-robustness}): after zero-forcing equalization, each fading realization presents the relay with an AWGN-like observation whose conditional noise variance is $\sigma^{2}/|h|^{2}$, i.e.\ a random effective SNR, and multi-SNR training already covers this mixture. Two honest qualifications apply. First, this adaptivity is an advantage over AF, whose fixed gain is tuned to an assumed SNR, but not over symbol-wise DF, whose $\operatorname{sign}(y)$ slicer is equally condition-agnostic; adaptivity therefore does not alter the H2 conclusion. Second, adaptation has been demonstrated only within the canonical Rayleigh ensemble; generalization to other channel families (Rician, bursty interference, phase noise, hardware non-linearities) remains untested and is identified as future work.

\section{Limitations}\label{sec:limitations}

Several limitations of this study should be acknowledged, as they define the boundary conditions under which the conclusions hold:

\begin{enumerate}
\def\labelenumi{\arabic{enumi}.}
\item
  \textbf{Modulation scope.} The canonical relay comparisons use QPSK, and the unknown-channel study (Chapter~\ref{sec:unknown-channels}) mainly uses BPSK, with a QPSK re-check of the unknown-ISI result (Section~\ref{sec:qpsk-unknown-channel}); 16-QAM enters only as a modulation-and-coding rung in the link-adaptation study, never as a relay-architecture question. In that sense higher-order constellations are not investigated: 16-QAM and denser grids break the per-axis relay formulation and require a joint $N$-class 2D relay whose output dimension scales with $M$, which is left to future work.
\item
  \textbf{Perfect CSI (canonical comparison).} The canonical comparison assumes perfect channel state information at each receiver; the finite-pilot and blind regimes of Chapter~\ref{sec:unknown-channels} relax this deliberately and are not limitations of that study. Within the canonical comparison, in practice channel estimation errors would affect the coherent compensation step and hence relay performance. Imperfect CSI introduces a model mismatch between the assumed and actual channel, which could differentially impact the relay strategies: AI relays might be more robust (having learned from noisy data) or less robust (having not been exposed to estimation errors during training).
\item
  \textbf{Static offline training.} AI relays are trained once on synthetic data and deployed without adaptation. While multi-SNR training covers the canonical fading ensemble (Section \ref{sec:channel-robustness}), and the unknown-channel study (H5) shows a fixed network can learn a \emph{stationary} unknown impairment, this protocol does not exploit online information; adaptive or continual training would be required if the deployment channel drifts over time.
\item
  \textbf{Single relay.} The framework considers a single relay node. Multi-relay cooperation introduces relay selection, power allocation, and cooperative diversity dimensions not captured in the two-hop model. With multiple relays, the system-level optimization (which relay to use, how to combine multiple relay observations) becomes as important as the per-relay processing function studied here.
\item
  \textbf{Two-hop only.} Extension to multi-hop networks with 3+ relays introduces noise accumulation (for AF) and error propagation (for DF) that grow with the number of hops. AI relays might show greater advantage in multi-hop settings where noise accumulation is more severe.
\item
  \textbf{Fixed window size.} The relay window size (\(w = 2\) or \(w = 5\)) is fixed and relatively small. Larger windows could provide more context for denoising, particularly in channels with temporal correlation (e.g., block fading). A systematic study of window-size optimization is deferred.
\item
  \textbf{No coding, in the core comparison.} The core comparison operates without error-correcting codes, so the DF baseline there is its symbol-level variant (Remark~\ref{rem:df-terminology}). Sections~\ref{sec:coded-block-df}--\ref{sec:coded-latency-capacity} measure a coded, practical finite-length coded block-DF baseline, its soft-decision counterparts, a rate-adaptive envelope over a modulation-and-coding grid, and that envelope re-derived under a round-trip latency budget, directly rather than leaving any of this as an assertion, but that study is confined to one code family, one frame length, and one channel model, and the relay and destination each decide once with no iterative exchange between them.
\item
  \textbf{Incomplete seed replication across architectures; unequal training budgets.} The original cross-architecture canonical comparison used one trained instance per architecture. Targeted follow-up studies since replicated the MLP size sweep, selected unknown-channel experiments, and equal-budget initialization sensitivity; a fully equalized multi-seed cross-architecture BER comparison remains incomplete. Where a result rests on a single trained instance, the confidence intervals accompanying it are Monte Carlo intervals over the \emph{evaluation} channel and bit realizations for that fixed model, not over the training procedure. Separately, the architectures do not all receive identical training budgets (e.g.\ the cGAN trains for 200 epochs against 100 for the other relays, Section~\ref{sec:relay-strategies}), a confound the normalized-parameter study (Section~\ref{sec:normalized-parameter-comparison}) does not remove, since it equalizes parameter count and input window but not epoch count. This gap has since been closed on one axis and remains open on the other. The MLP size axis was re-run across three independent initializations per configuration (Section~\ref{sec:minimum-relay-size}), and the resulting across-initialization spread is reported in Table~\ref{tbl:seed-spread}: median $0.180$~dB over 152 configurations, but strongly channel-dependent, from $0.008$~dB on the canonical channel to $1.231$~dB on the nonlinear-bias channel. That spread is itself the reason the single-run caveat mattered, since on four of the nine channels it exceeds the entire effect of varying model size. The \emph{cross-architecture} comparison, however, is still a single trained instance per architecture with unequal training budgets, so the claim that architecture is secondary to capacity (Section~\ref{sec:parameter-normalization-complexity-trade-off}) remains an observation from one run rather than a seed-robust result; a multi-seed re-run of that comparison with equalized budgets is the direct way to close what is left and is not performed in this thesis.
\end{enumerate}

\section{Future Work}\label{sec:future-work}

Several directions warrant further investigation:

\begin{enumerate}
\def\labelenumi{\arabic{enumi}.}
\item
  \textbf{Time-selective fading.} The present study addresses channels that are static within a block. Extending the learned relay to time-selective fading with continuously varying unknown phase (e.g.\ Jakes/Clarke Doppler) is a natural next step, likely requiring explicitly phase-aware architectures (complex-valued state, learned phase unwrapping, or a training objective aligned with noncoherent detection rather than symbol MSE) evaluated against prediction-based noncoherent receivers.
\item
  \textbf{Higher-order constellations (16-QAM and denser).} Constellations with more than two levels per axis break the per-axis relay formulation and need a joint $N$-class 2D relay whose output dimension grows with $M$; that reformulation, and whether the complexity ordering discussed under H3 shifts with constellation density, are open questions this thesis does not address.
\item
  \textbf{Imperfect CSI.} Introduce channel estimation errors to assess robustness of AI relay processing under realistic conditions.
\item
  \textbf{A BER-optimal (BCJR/APP) benchmark for the unknown-ISI channel.} Section~\ref{sec:unknown-channel-experiment} benchmarks against genie-CSI Viterbi MLSE, the optimal \emph{sequence} detector. Sequence-ML and bit-BER-optimal detection are distinct criteria on any channel with memory, at both modulation orders studied here, so the benchmark is not BER-optimal in either case; for Gray-coded QPSK there is the additional effect that a wrong trellis path can differ from the truth by one or two bits with unequal probability. That second effect has since been measured (Section~\ref{sec:qpsk-unknown-channel}) and is small here: the learned relay and the trellis lose $1.073$ and $1.090$ bits per symbol error respectively, so the Gray map moves the two detectors almost equally. A BCJR/APP detector run over the same 3-tap trellis, which directly minimizes per-bit error probability rather than sequence error probability, would test whether the classical benchmark's margin over the learned relay changes once it is BER-optimal rather than sequence-optimal. It is no longer needed to explain any observation in this thesis --- the QPSK ordering that once appeared to require one turned out to rest on a benchmark that had not been given the channel's fading gains --- but the margin itself has not been measured against a BER-optimal detector at either modulation order, and that remains open.
\item
  \textbf{Online learning.} Develop relay strategies that adapt their parameters during operation, tracking time-varying channel conditions.
\item
  \textbf{Multi-relay networks.} Extend to cooperative relay selection and multi-hop routing with AI-optimized relays at each node.
\item
  \textbf{Other channels and MIMO.} Evaluate the canonical conclusions on additional channel families (Rician line-of-sight fading, and AWGN treated as a full operating point rather than the bounded calibration-plus-companion role it has here) and on a multi-antenna second hop, where destination equalization interacts with relay processing and the additivity of the two gains becomes a testable question.
\item
  \textbf{End-to-end learning.} Train the entire communication chain (modulation, relay, demodulation) jointly using autoencoder-based approaches, and compare against the modular classical-plus-learned-relay cascade studied here.
\item
  \textbf{Coded block-DF baseline (measured; open items remain).} Remark~\ref{rem:df-terminology} scopes the core comparison's DF baseline to the uncoded, symbol-level variant. Section~\ref{sec:coded-block-df} measures the coded, practical finite-length coded block-DF baseline this item once only proposed, together with the soft-information relays (BCJR at the relay, posterior-mean read-out of the learned relay) and an adaptive-rate envelope over a modulation-and-coding grid. The substantive finding is mechanistic: a relay that decodes the code has spent its redundancy before transmitting, so a wrong decode arrives as a valid-but-wrong codeword the destination cannot repair, whereas a relay that only denoises leaves the redundancy intact. Rate selection then dominates the relay-decoding question by an order of magnitude, unless a latency deadline is imposed, at which point block-DF's double buffering reopens the gap. What remains open: one code family, one frame length, one channel model; the soft block-DF relay was given the nominal noise variance rather than per-symbol CSI, which caused its 20~dB mis-calibration; iterative (turbo-style) relay--destination exchange was not evaluated; and the latency analysis uses structural buffer counts rather than a $10^{-5}$-class reliability target, so it does not establish ultra-reliable low-latency performance.
\item
  \textbf{Longer relay windows with state-space architectures.} Future work should explore whether longer relay windows (e.g., 128--512 symbols) combined with a chunk-parallel state-space architecture can improve both BER performance and processing speed simultaneously; the training-time crossover between sequential and chunk-parallel formulations, removed from this thesis's scope, is the natural companion measurement.
\end{enumerate}

\section{Hypotheses Revisited}\label{sec:conclusions}

Chapter~\ref{sec:summary} states what this thesis found, and Table~\ref{tbl:table33} there records the verdict on each hypothesis. This section addresses the different question of \emph{how} each hypothesis fared under measurement, and why three of the five required qualification.

\paragraph{Do the other hypotheses survive the size sweep?} The sweep of Section~\ref{sec:minimum-relay-size} re-measures four of the five hypotheses' subject matter with a different relay class, a corrected two-hop model, per-channel classical comparators and three initializations per configuration, so it is worth asking explicitly whether any of them moves. Only H3 does.

H1 is reconfirmed on the canonical channel: the smallest matching relay reaches $0.3294$ at 0~dB and $0.2198$ at 4~dB against AF's $0.4088$ and $0.3140$, so the low-SNR advantage over AF is not an artifact of the particular architectures originally compared. H2 is likewise reconfirmed and, for the first time, quantified in the units that matter: on the canonical channel the learned relay trails symbol-wise DF by $0.05$ to $0.09$~dB, which is a real margin in DF's favour but two orders of magnitude smaller than the ``$+2.06\%$ relative BER'' figure that the same measurements produce at 20~dB. H2's direction stands; its magnitude was previously overstated by the choice of metric.

H5 receives independent corroboration rather than merely surviving. Chapter~\ref{sec:unknown-channels} reports the learned relay $1$--$1.5$~dB behind genie-CSI Viterbi on the three-tap ISI channel, using a window of 11. The sweep, which caps the window at 7 and uses a differently trained relay, measures $+1.67$~dB at that width, with a clean progression of $+12.95$, $+3.60$, $+2.75$ and $+1.67$~dB at windows 1, 3, 5 and 7. Extrapolating that trend to the window of 11 actually used in Chapter~\ref{sec:unknown-channels} gives $+0.79$~dB, bracketing the reported figure from below as the measured points bracket it from above. Two independent measurements, differing in relay class, training procedure, two-hop model and Monte Carlo budget, therefore agree on the size of the gap to the model-aware detector, and the residual difference between them is explained by window width alone --- which is the same variable the sweep identifies as binding.

H4 needs qualifying, and the qualification could not be inferred from the size sweep. That sweep measures initialization variance on MLPs, where the canonical channel is the most reproducible case anywhere in this thesis at $0.008$~dB median, and an earlier version of this section offered that as reassurance for H4's single-seed equal-budget comparison while noting it was not proof, the sequence models never having been measured. Measuring them directly reverses the conclusion for one of the six.

Each equal-budget architecture was retrained from three seeded initializations and evaluated on identical channel draws and payload bits. The resulting spread in SNR penalty against DF is tabulated as Table~\ref{tbl:seed-spread-3k} in Appendix~\ref{sec:appendix-f-minimum-size}.

Four of the five are stable to within $0.16$~dB and the two Mamba variants are the most reproducible architectures measured, at $0.019$ and $0.041$~dB. The Transformer is not: its three initializations span $0.906$~dB, five times the effect H4's high-SNR claim must resolve, with one run at $+2.039$~dB against two at approximately $+1.14$~dB. The instability is therefore specific to that architecture rather than generic to sequence models, which is the opposite of what a capacity or inductive-bias account would predict.

\paragraph{What the Transformer's instability is.} Running eight initializations rather than three, and recording the training objective alongside the BER, identifies the cause as an outright optimization failure rather than a property of the architecture's solution. Figures~\ref{fig:transformer-seed-curves} and~\ref{fig:transformer-loss-penalty} show the eight runs and the loss that separates them. Seven of the eight seeds converge to a final training loss between $0.0018$ and $0.0023$ and land within $0.033$~dB of one another, which is tighter than the Mamba control's $0.041$~dB. The eighth converges to $0.0256$, eleven times the median, and is the $+2.039$~dB run. Final training loss and SNR penalty correlate at $+0.998$ across the eight runs, but that figure should not be read as a general relationship between loss and error rate: it is produced entirely by the single outlier's leverage. Among the seven runs that converge, loss varies by roughly $30\%$ with no detectable effect on BER ($r=-0.38$, $p=0.41$). The correct reading is a two-cluster one --- the failed run is worse on both measures simultaneously, which is what identifies it as a failure of optimization rather than an alternative solution reached by a different route, while within the converged cluster the training loss carries no information about which relay will perform better. Training loss is therefore a reliable detector of catastrophic non-convergence and useless as a fine-grained quality signal, and nothing here suggests the objective is mismatched to the error rate the relay is judged on.

Three practical points follow. The failure occurs in roughly one initialization in eight at this budget, so a three-seed sample had about a one-in-three chance of containing one and the published single-seed run had one in eight of being it. It is trivially detectable at training time, the bad run's loss being an order of magnitude off, so seed selection on final training loss would remove the instability entirely at no evaluation cost. And it is \emph{not} detectable from validation accuracy, which reads $0.9971$ for the failed seed and spans $0.9971$ to $0.9992$ across the healthy ones: the metric the training loop already monitors does not separate the failure, which is presumably why it was never noticed. Excluding the single failed run, the Transformer is as reproducible as any other architecture measured, so the correct statement is not that the Transformer is an unstable architecture but that it fails to converge occasionally and that the failure is silent under the diagnostics in use.

Two consequences follow for H4. H4's low-SNR finding stands: the feedforward and sequence groups separate by roughly a decibel here ($+0.15$ against $+1.1$~dB), far outside any measured spread, so the $17.4\%$ split at 0~dB is not a seeding artifact. H4's high-SNR convergence claim does not stand as a seed-robust result, because the $4.1\%$ spread it rests on is smaller than the Transformer's own run-to-run variation; it should be read as an observation from one run of each architecture. Notably, this was not detectable from the published experiments: none of the three sequence checkpoints seeds \texttt{torch} at all, so the reported comparison was drawn from an uncontrolled initialization and its variance was not merely unreported but unmeasurable after the fact.

The main conclusions, in order of significance, are:

These six findings are enumerated in Chapter~\ref{sec:summary} and are not restated here; what follows concerns the equal-budget comparison behind the fourth of them.

It is important to note that sequence-detection methods such as Viterbi MLSE appear only in the unknown-channel study of Chapter~\ref{sec:unknown-channels}, where channel memory is introduced explicitly through a 3-tap FIR/ISI channel. The canonical system model of Chapters~\ref{sec:introduction}--\ref{sec:experiments} remains memoryless and therefore requires only per-symbol processing.

The unknown-channel contribution (Chapter~\ref{sec:unknown-channels}) settles the complementary question of \emph{when the learned relay is not merely competitive but necessary}:

The flat-channel control isolates the cause and is the part worth restating here: against unknown phase, gain and per-branch asymmetry, all memoryless, the classical relay does not fail and the learned relay only matches it (largest difference $0.0012$ over the tabulated SNRs). Parameter uncertainty alone is therefore not what learning mitigates; failure appears only with unmodeled \emph{memory}. The remaining findings of that chapter --- the analytic $0.25$ floor the memoryless relays sit at, the margin to pilot-aided MLSE, the behaviour with no pilots at all, and the pilot budget at which the crossover occurs --- are enumerated in Chapter~\ref{sec:summary} and established in Chapter~\ref{sec:unknown-channels}; they are not repeated here.

Taken together these results place neural relay processing as a complement to classical methods rather than a replacement for them, and Chapter~\ref{sec:summary} states the resulting deployment recommendation and the two boundaries that bound it. The interpretive point this chapter adds is the one the individual findings do not state on their own: \textbf{model complexity should be matched to task complexity, and the task's complexity is set by the channel rather than by the constellation or the budget}. Where the channel is memoryless the relay denoising task is small enough that minimal architectures suffice and the choice between feedforward and sequential paradigms matters less than sizing; where the channel has memory, the same reasoning raises the floor rather than lowering it, and does so through the relay's window rather than through its parameter count.

\begin{center}\rule{0.5\linewidth}{0.5pt}\end{center}
